\section{Hard Instances and Proof of Theorem~\ref{thm:lower-bound}}
\label{sec:lb-proof}

In Section~\ref{sec:lower-bound}, we established a lower bound on the iteration
complexity of the SPLALM family that holds for every choice of constant center
stepsizes. In this section, we prove this result by constructing hard instances.

Fix \(\epsilon\) and \(\theta\) satisfying the conditions of
Theorem~\ref{thm:lower-bound}. The lower bound has the quantified form
\[
\forall\,(\alpha,\beta)\in[0,1]^2,\quad
\exists\,I\in\mathcal I(\mathbf C_0,\Delta):\quad
T^{\mathrm{KKT}}_\epsilon(\theta,\alpha,\beta;I)
\geq \Omega(\Delta\epsilon^{-4}).
\]
Thus, different choices of \((\alpha,\beta)\) may require different hard
instances. For each pair, we construct an instance in
\(\mathcal I(\mathbf C_0,\Delta)\) on which
\(\mathrm{SPLALM}_{\theta,\alpha,\beta}\), started from the zero
initialization prescribed in Definition~\ref{def:splalm}, requires at least
the stated number of iterations to produce an \(\epsilon\)-KKT output.
For the outputs of \(\mathrm{SPLALM}_{\theta,\alpha,\beta}\), we write
\(R_k:=\mathcal R_{\mathrm{KKT}}(\bx^{k+1},\by^{k+1},\bq^k)\) for \(k\ge0\).

We split \([0,1]^2\) into three parts by the threshold
\[
\alpha_\epsilon:=1024\max\{1,4\tau,4\gamma\tau\}\,\epsilon^2,
\]
namely
\[
\begin{aligned}
[0,1]^2
={}&\bigl\{(\alpha,\beta)\in[0,1]^2:\alpha\le\alpha_\epsilon\bigr\}\\
&\cup\bigl\{(\alpha,\beta)\in[0,1]^2:\alpha>\alpha_\epsilon,\ \beta>0\bigr\}\\
&\cup\bigl\{(\alpha,\beta)\in[0,1]^2:\alpha>\alpha_\epsilon,\ \beta=0\bigr\},
\end{aligned}
\]
and use a separate hard instance for each part. For
\(\alpha\le\alpha_\epsilon\), the bound obtained in
Section~\ref{subsec:lb-plateau} is of order \(\Delta\epsilon^{-4}\). This
threshold ensures the parameter estimates required by the rotating instance
in Section~\ref{subsec:lb-rotation}.

\subsection{Reduction to a gradient descent--ascent recursion}
\label{subsec:gda-embedding}

Where \(\by(\bx,\blmbd)=\zero\), the updates of
\(\mathrm{SPLALM}_{\theta,\alpha,\beta}\) reduce to gradient
descent--ascent steps on a smooth saddle function. To state this, define
\[
\mathcal R_\rho:=\bigl\{(\bx,\blmbd):\ \blmbd+\rho\calA(\bx)\in\partial h(\zero)\bigr\},
\qquad
\Psi_\rho(\bx,\blmbd):=f(\bx)+\frac{\rho}{2}\|\calA(\bx)\|^2+\langle\blmbd,\calA(\bx)\rangle .
\]

\begin{proposition}[Reduction on \(\mathcal R_\rho\)]
\label{prop:reduction}
Let \(\rho>0\) and \((\bx,\blmbd)\in\mathcal R_\rho\). Then
\(\by(\bx,\blmbd)=\zero\), \(\bq(\bx,\blmbd)=\blmbd+\rho\calA(\bx)\), and
\[
\nabla f(\bx)+J_{\calA}(\bx)^\top\bq(\bx,\blmbd)=\nabla_{\bx}\Psi_\rho(\bx,\blmbd),
\qquad
\calA(\bx)-\by(\bx,\blmbd)=\nabla_{\blmbd}\Psi_\rho(\bx,\blmbd).
\]
Moreover,
\(\calD_\rho(\bx,\blmbd)=\Psi_\rho(\bx,\blmbd)+h(\zero)\).
\end{proposition}

\begin{proof}
Set \(\bv:=\calA(\bx)+\blmbd/\rho\). By~\eqref{eq:prox-multiplier}, \(\by(\bx,\blmbd)\)
is the unique minimizer of \(h(\by)+\frac{\rho}{2}\|\by-\bv\|^2\), and
\(\by\) is this minimizer if and only if \(\rho(\bv-\by)\in\partial h(\by)\).
For \(\by=\zero\) the inclusion reads
\(\blmbd+\rho\calA(\bx)\in\partial h(\zero)\), which holds because
\((\bx,\blmbd)\in\mathcal R_\rho\). Hence \(\by(\bx,\blmbd)=\zero\) and
\(\bq(\bx,\blmbd)=\blmbd+\rho\calA(\bx)\). Since
\(\nabla_{\bx}\bigl(\frac{\rho}{2}\|\calA(\bx)\|^2\bigr)=\rho J_{\calA}(\bx)^\top\calA(\bx)\),
\[
\nabla f(\bx)+J_{\calA}(\bx)^\top\bigl(\blmbd+\rho\calA(\bx)\bigr)
=\nabla_{\bx}\Psi_\rho(\bx,\blmbd),
\]
and \(\calA(\bx)-\zero=\nabla_{\blmbd}\Psi_\rho(\bx,\blmbd)\).
Finally, completing the square in~\eqref{eq:reduced-al} gives
\(\calD_\rho(\bx,\blmbd)=f(\bx)+\min_{\by}\{h(\by)+\frac{\rho}{2}\|\by-\bv\|^2\}
-\frac{1}{2\rho}\|\blmbd\|^2\), and the minimum equals
\(h(\zero)+\frac{\rho}{2}\|\bv\|^2\) because its minimizer is \(\zero\). Hence
\[
\calD_\rho(\bx,\blmbd)
=f(\bx)+h(\zero)+\frac{\rho}{2}\|\calA(\bx)\|^2+\langle\blmbd,\calA(\bx)\rangle
+\frac{1}{2\rho}\|\blmbd\|^2-\frac{1}{2\rho}\|\blmbd\|^2
=\Psi_\rho(\bx,\blmbd)+h(\zero).
\]
\end{proof}

For \(\theta=(\rho,s,\gamma,\eta,\tau)\) and \((\alpha,\beta)\in[0,1]^2\), consider
the recursion
\begin{equation}\label{eq:gda-recursion}
\begin{aligned}
\bx^{k+1}&=\bx^k-\eta\Bigl[\nabla_{\bx}\Psi_\rho(\bx^k,\blmbd^k)
  +\frac{1}{\gamma}(\bx^k-\bz^k)\Bigr],\\
\blmbd^{k+1}&=\blmbd^k+\tau\Bigl[\nabla_{\blmbd}\Psi_\rho(\bx^{k+1},\blmbd^k)
  -\frac{1}{s}(\blmbd^k-\bmu^k)\Bigr],\\
\bz^{k+1}&=\bz^k+\alpha(\bx^{k+1}-\bz^k),\qquad
\bmu^{k+1}=\bmu^k+\beta(\blmbd^{k+1}-\bmu^k).
\end{aligned}
\end{equation}

\begin{proposition}[Embedding]
\label{prop:embedding}
Let \(\rho>0\), and let the sequence
\(\bw^k:=(\bx^k,\blmbd^k,\bz^k,\bmu^k)\) be generated by
\(\mathrm{SPLALM}_{\theta,\alpha,\beta}\). Suppose that for some \(K\ge1\),
\((\bx^k,\blmbd^k)\in\mathcal R_\rho\) and
\((\bx^{k+1},\blmbd^k)\in\mathcal R_\rho\) for all \(0\le k<K\).
Then \eqref{eq:gda-recursion} holds for \(0\le k<K\), and
\[
R_k=\max\bigl\{\|\nabla_{\bx}\Psi_\rho(\bx^{k+1},\blmbd^k)\|,\;
\|\nabla_{\blmbd}\Psi_\rho(\bx^{k+1},\blmbd^k)\|\bigr\},
\qquad 0\le k<K.
\]
\end{proposition}

\begin{proof}
Apply Proposition~\ref{prop:reduction} at \((\bx^k,\blmbd^k)\) in the
\(\bx\)-update of~\eqref{eq:splalm-updates}, and at \((\bx^{k+1},\blmbd^k)\) in
the \(\blmbd\)-update. The same point gives \(\by^{k+1}=\zero\) and
\(\bq^k=\blmbd^k+\rho\calA(\bx^{k+1})\), so the formula for \(R_k\) follows
from~\eqref{eq:kkt-residual}.
\end{proof}

Recursion~\eqref{eq:gda-recursion} is the update of
DS-GDA~\cite{zheng2023universal} for the saddle function \(\Psi_\rho\), with
the same order of updates, whenever the projections in DS-GDA are inactive.
Proposition~\ref{prop:embedding} identifies
\(\mathrm{SPLALM}_{\theta,\alpha,\beta}\) with this recursion only while both
evaluation points \((\bx^k,\blmbd^k)\) and \((\bx^{k+1},\blmbd^k)\) lie in
\(\mathcal R_\rho\); outside \(\mathcal R_\rho\), the two recursions differ.

\subsection{Small primal center stepsize $\alpha$}
\label{subsec:lb-plateau}

In this subsection, we treat the pairs \((\alpha,\beta)\) with
\(\alpha\le\alpha_\epsilon\). We construct a hard instance on which, for
every \(\beta\in[0,1]\), \(\mathrm{SPLALM}_{\theta,\alpha,\beta}\) requires
more than \(\Delta/(16\gamma\alpha\epsilon^2)\) iterations to produce an
\(\epsilon\)-KKT output. The hard instance is given below.

\begin{example}
\label{ex:plateau}
Let \(0<\epsilon\le\min\{1/2,\sqrt{\Delta}/2\}\) and
\(x_\star:=\Delta/(2\epsilon)+\epsilon\). Define \(\calA\equiv0\),
\(h(y)=|y|\), and \(f:\RR\to\RR\) by
\[
f(x)=
\begin{cases}
2\epsilon(x_\star-x)-2\epsilon^2, & x<x_\star-2\epsilon,\\[4pt]
\dfrac12(x-x_\star)^2, & |x-x_\star|\le2\epsilon,\\[4pt]
2\epsilon(x-x_\star)-2\epsilon^2, & x>x_\star+2\epsilon.
\end{cases}
\]
Then \(f\) is \(1\)-smooth with \(\|\nabla f\|\le2\epsilon\le1\), \(\calA\) satisfies
\(L_{\calA}=G_{\calA}=0\), and \(h\) is convex and \(1\)-Lipschitz. Since
\(\epsilon\le\sqrt{\Delta}/2\), we have \(x_\star\ge2\epsilon\), thus \(f(0)=2\epsilon x_\star-2\epsilon^2=\Delta\). Moreover, \(\inf f=f(x_\star)=0\). Hence
\((f,\calA,h)\in\mathcal I(\mathbf C_0,\Delta)\).
\end{example}

\begin{remark}
\label{rem:plateau}
The construction ensures that \(\inf f=0\) and \(f(0)=\Delta\), so the
initial objective gap is exactly \(\Delta\). This allows us to make the
dependence of the iteration lower bound on \(\Delta\) explicit.
\end{remark}

This instance shows that a small primal center stepsize can force
\(\Omega\bigl(\Delta/(\gamma\alpha\epsilon^2)\bigr)\) iterations even on a
smooth one-dimensional problem, regardless of the dual center stepsize.
The iterates must traverse a long constant-gradient interval before an
output can satisfy the KKT tolerance. A small \(\alpha\) limits how quickly
the center moves, and the primal variable stays within a bounded distance
of that center. Since the multipliers remain zero, every output on this
interval has KKT residual \(2\epsilon>\epsilon\). The next lemma bounds the
distance traveled and converts it into a lower bound on the number of
iterations.

\begin{lemma}
\label{lem:plateau}
Let \((\alpha,\beta)\in[0,1]^2\), and suppose that \(\epsilon\) satisfies the
condition in Example~\ref{ex:plateau} and \(\epsilon^2\le\Delta/(16\gamma)\).
Run \(\mathrm{SPLALM}_{\theta,\alpha,\beta}\) on Example~\ref{ex:plateau}.
Then the following hold.
\begin{enumerate}[label=(\roman*)]
\item For all \(k\ge0\), \(\lambda^k=\mu^k=0\), \(y^{k+1}=0\), and
  \(R_k=\|\nabla f(x^{k+1})\|\).
\item If \(\alpha>0\), then for every integer \(k\) with \(0\le k\le\Delta/(16\gamma\alpha\epsilon^2)\),
  \[
  \|z^k-z^0\|\le2\alpha\gamma\epsilon\,k,\qquad
  \|x^k-z^k\|\le2\gamma\epsilon,
  \]
  and
  \[
  \|x^k-x^0\|\le2\gamma\epsilon(1+\alpha k)
  \le x_\star-2\epsilon .
  \]
  When \(\alpha=0\), these bounds hold for all \(k\ge0\).
\item Let \(I\) denote the instance of Example~\ref{ex:plateau}. Then
  \(T^{\mathrm{KKT}}_\epsilon(\theta,\alpha,\beta;I)>\Delta/(16\gamma\alpha\epsilon^2)\)
  if \(\alpha>0\), and \(T^{\mathrm{KKT}}_\epsilon(\theta,\alpha,\beta;I)=+\infty\)
  if \(\alpha=0\).
\end{enumerate}
\end{lemma}

\begin{proof}
(i) Since \(\calA\equiv0\) and \(\partial h(0)=[-1,1]\), every pair
\((x,0)\) lies in \(\mathcal R_\rho\), and \(\Psi_\rho(x,\lambda)=f(x)\). The
\(\lambda\)-update in~\eqref{eq:gda-recursion} reads
\(\lambda^{k+1}=\lambda^k-\frac{\tau}{s}(\lambda^k-\mu^k)\). Starting from
\(\lambda^0=\mu^0=0\), induction gives \(\lambda^k=\mu^k=0\) for all \(k\).
Proposition~\ref{prop:embedding} then gives \(y^{k+1}=0\) and
\(R_k=\max\{\|\nabla f(x^{k+1})\|,0\}\).

(ii) We use the Euclidean norm, which coincides with the absolute value
in one dimension. By (i), the primal update gives
\[
x^{k+1}-z^k
=\left(1-\frac{\eta}{\gamma}\right)(x^k-z^k)-\eta\nabla f(x^k).
\]
Since \(\|\nabla f(x)\|\le2\epsilon\) for every \(x\) and
\(0<\eta\le\gamma\),
\[
\|x^{k+1}-z^k\|
\le\left(1-\frac{\eta}{\gamma}\right)\|x^k-z^k\|+2\eta\epsilon.
\]
The center update implies
\[
\|x^{k+1}-z^{k+1}\|=(1-\alpha)\|x^{k+1}-z^k\|,
\qquad
\|z^{k+1}-z^k\|=\alpha\|x^{k+1}-z^k\|.
\]
Starting from \(x^0=z^0\), induction gives
\(\|x^k-z^k\|\le2\gamma\epsilon\) and
\(\|x^{k+1}-z^k\|\le2\gamma\epsilon\) for all \(k\ge0\).
Consequently,
\[
\|z^k-z^0\|
\le\sum_{j=0}^{k-1}\|z^{j+1}-z^j\|
\le2\alpha\gamma\epsilon k.
\]
Using \(x^0=z^0=0\), we obtain
\[
\|x^k-x^0\|
\le\|x^k-z^k\|+\|z^k-z^0\|
\le2\gamma\epsilon(1+\alpha k).
\]
If \(\alpha>0\), then for every integer \(k\le\Delta/(16\gamma\alpha\epsilon^2)\), the
assumptions \(\epsilon^2\le\Delta/(16\gamma)\) and
\(\epsilon^2\le\Delta/4\) imply
\[
\|x^k-x^0\|
\le2\gamma\epsilon+\frac{\Delta}{8\epsilon}
\le\frac{\Delta}{4\epsilon}
\le\frac{\Delta}{2\epsilon}-\epsilon
=x_\star-2\epsilon.
\]
In particular, \(x^k\le\|x^k-x^0\|\le x_\star-2\epsilon\), so the
iterate lies in the constant-gradient region. When \(\alpha=0\), the
center remains at \(z^k=z^0=0\), and the same conclusion holds for every
\(k\ge0\).

(iii) Suppose \(\alpha>0\). By (ii), every output with
\(k+1\le\Delta/(16\gamma\alpha\epsilon^2)\) satisfies
\(x^{k+1}\le x_\star-2\epsilon\). By (i), its residual is
\(R_k=\|\nabla f(x^{k+1})\|=2\epsilon>\epsilon\).
Thus, none of the first
\(\lfloor\Delta/(16\gamma\alpha\epsilon^2)\rfloor\) outputs is an
\(\epsilon\)-KKT point. Definition~\ref{def:complexity-measure} yields
\[
T^{\mathrm{KKT}}_\epsilon(\theta,\alpha,\beta;I)
\ge\left\lfloor\frac{\Delta}{16\gamma\alpha\epsilon^2}\right\rfloor+1
>\frac{\Delta}{16\gamma\alpha\epsilon^2}.
\]
When \(\alpha=0\), (ii) keeps every iterate in the constant-gradient region.
Hence \(R_k=2\epsilon>\epsilon\) for all \(k\ge0\), and
\(T^{\mathrm{KKT}}_\epsilon(\theta,\alpha,\beta;I)=+\infty\).
\end{proof}

For \(0<\alpha\le\alpha_\epsilon\), Lemma~\ref{lem:plateau}~(iii) gives
\[
T^{\mathrm{KKT}}_\epsilon(\theta,\alpha,\beta;I)
>\frac{\Delta}{16\gamma\alpha\epsilon^2}
\ge\frac{\Delta\,\epsilon^{-4}}{16384\,\gamma\max\{1,4\tau,4\gamma\tau\}}.
\]

\subsection{Large primal center stepsize \texorpdfstring{$\alpha$}{alpha}, and \texorpdfstring{$\beta>0$}{beta>0}}
\label{subsec:lb-rotation}

In this subsection, we prove the lower bound for
\(\alpha>\alpha_\epsilon\) and \(\beta>0\). For every such pair, we construct
a hard instance on which SPLALM, initialized at zero, never produces an
\(\epsilon\)-KKT point.

% On the region where the objective is quadratic, and as long as the prox
% outputs vanish, every update of SPLALM is linear. One iteration is therefore
% a fixed matrix acting on the current state, and the eigenvalues of this
% matrix determine the trajectory. If all eigenvalues had modulus less than
% one, the iterates would converge to a KKT point and the residual would
% vanish; if some eigenvalue had modulus greater than one, the iterates would
% leave the quadratic region. We therefore choose the parameters so that the
% matrix has a non-real eigenvalue of modulus one. The corresponding component
% of the trajectory rotates with constant amplitude and never settles at a
% fixed point. The initialization is fixed at zero and is not on this
% component. We show that the trajectory from zero differs from the rotation
% by a decaying term that stays small at every iteration, so the residual
% stays above \(\epsilon\). The hard instance is given below.

Motivated by the quadratic hard instances in~\cite[Example~5.1]{li2026smoothing}
and~\cite[Appendix~C.1]{mahdavinia2022tight}, we construct the following instance.

\begin{example}
\label{ex:rotation}
Let \(b,\kappa,X>0\) and \(x_c,\lambda_c\in\mathbb C\) be given in
Appendix~\ref{app:lb-rotation-params}, and set \(R_*:=\frac54X\). Define
\(\calA(x)=b(x-x_c)\), \(h(y)=\|y\|_2\), and
\[
f(x)=\psi\bigl(\|x-x_c\|\bigr)-b\langle\lambda_c,x-x_c\rangle,
\qquad
\psi(r)=
\begin{cases}
-\dfrac{\kappa}{2}r^2, & 0\le r\le R_*,\\[4pt]
\dfrac{\kappa}{2}(r-2R_*)^2-\kappa R_*^2, & R_*\le r\le3R_*,\\[4pt]
\kappa R_*(r-3R_*)-\dfrac{\kappa}{2}R_*^2, & r\ge3R_*.
\end{cases}
\]
\end{example}
To establish this lower bound, we analyze the SPLALM updates on
Example~\ref{ex:rotation}. Specifically, since \(h=\|\cdot\|_2\), Proposition~\ref{prop:reduction} ensures that both proximal outputs vanish
whenever
\[
\max_{j\in\{k,k+1\}}\|\lambda^k+\rho\calA(x^j)\|\le1.
\] 
Under this condition, SPLALM follows~\eqref{eq:gda-recursion}. In the regime where \(\|x-x_c\|\le R_*\) and \(\|\lambda+\rho\calA(x)\|\le1\),
the reduced function takes the quadratic form
\[
\Psi_\rho(x,\lambda)
=-\frac{a}{2}\|x-x_c\|^2+b\langle\lambda-\lambda_c,x-x_c\rangle .
\]
Here \(a:=\kappa-\rho b^2\). The gradients of this quadratic function are
linear in \(x-x_c\) and \(\lambda-\lambda_c\). 
Therefore, let define the iteration sequence \(\bw^k:=(x^k,\lambda^k,z^k,\mu^k)\) and
\(\bw_c:=(x_c,\lambda_c,x_c,\lambda_c)\). Whenever the proximal condition
above and \(\|x^k-x_c\|\le R_*\) hold, the four updates take the linear form
\[
\bw^{k+1}-\bw_c=\mathbb T(\bw^k-\bw_c),
\] with $\mathbb T$ defined as 
\begin{equation}
\label{eq:rot-matrix}
\mathbb T=\begin{pmatrix}
1+\eta a-\tfrac{\eta}{\gamma} & -\eta b & \tfrac{\eta}{\gamma} & 0\\[2pt]
\tau b\bigl(1+\eta a-\tfrac{\eta}{\gamma}\bigr) & 1-\tfrac{\tau}{s}-\eta\tau b^2
  & \tfrac{\eta\tau b}{\gamma} & \tfrac{\tau}{s}\\[2pt]
\alpha\bigl(1+\eta a-\tfrac{\eta}{\gamma}\bigr) & -\alpha\eta b
  & 1-\alpha+\tfrac{\alpha\eta}{\gamma} & 0\\[2pt]
\beta\tau b\bigl(1+\eta a-\tfrac{\eta}{\gamma}\bigr)
  & \beta\bigl(1-\tfrac{\tau}{s}-\eta\tau b^2\bigr)
  & \tfrac{\beta\eta\tau b}{\gamma} & 1-\beta+\tfrac{\beta\tau}{s}
\end{pmatrix}.
\end{equation}
Under the conditions above, analyzing the SPLALM iterates reduces to
studying the spectral properties of \(\mathbb T\).
The core idea is to construct a bounded trajectory whose KKT residual remains
above \(\epsilon\). We choose \(a\) and \(b\) so that \(\mathbb T\) has a
non-real eigenvalue of modulus one, which gives a component that does not
decay.

The construction in~\cite[Section~5.1.1]{li2026smoothing} chooses both the
hard instance and the initial point, with the latter lying along an
eigenvector of the iteration matrix. In our setting, the initial point is
fixed at zero, and we choose the instance alone. We therefore need to
obtain the required trajectory from this prescribed initial point. To this
end, we choose the centers \(x_c\) and \(\lambda_c\) in the instance so that
the shifted initial state is the sum of an eigenvector of \(\mathbb T\)
and a decaying component. The following lemma gives the spectral
properties and this decomposition for the parameters defined
in Appendix~\ref{app:lb-rotation-params}.

\begin{lemma}
\label{lem:rot-spectral}
Under the conditions of Lemma~\ref{lem:rot-estimates}, let
\(\mathbb T\) be the matrix in~\eqref{eq:rot-matrix}. Then:
\begin{enumerate}[label=(\roman*)]
\item \(\mathbb T\bv=\zeta\bv\), where \(|\zeta|=1\) and \(\zeta\notin\RR\);
\item \(\br\) lies in a \(\mathbb T\)-invariant subspace on which the spectral
  radius of \(\mathbb T\) is less than one;
\item \(\bv+\br=-\bw_c\).
\end{enumerate}
\end{lemma}
 
By Lemma~\ref{lem:rot-spectral}, the linear recursion initialized at zero
has the explicit solution
\[
\bw^k-\bw_c=\zeta^k\bv+\mathbb T^k\br.
\]
We next bound \(\mathbb T^k\br\) uniformly in \(k\) to verify that this
sequence stays in the quadratic region, both proximal outputs vanish,
and the KKT residual remains above \(\epsilon\).

\begin{lemma}
\label{lem:rot-estimates}
Let \(\alpha>\alpha_\epsilon\), \(0<\beta\le1\), \(0<\epsilon<c_2\), and
\(\epsilon(\rho+s)\le c_1\), with \(c_1,c_2\) as specified in
Theorem~\ref{thm:lower-bound}. Then:
\begin{enumerate}[label=(\roman*)]
\item the instance of Example~\ref{ex:rotation} belongs to
  \(\mathcal I(\mathbf C_0,\Delta)\);
\item for all \(k\ge0\),
\[
\|(\mathbb T^k\br)_x\|\le\frac X4,
\qquad
\|(\mathbb T^k\br)_\lambda\|\le
s\Bigl(\frac{\gamma}{\eta}+\frac54\Bigr)bX.
\]
\end{enumerate}
\end{lemma}
The proofs of Lemmas~\ref{lem:rot-spectral} and~\ref{lem:rot-estimates} are
given in Appendix~\ref{app:lb-rotation}.

Using these estimates, we now show that the sequence generated by the linear
recursion coincides with the SPLALM iterates and that its KKT residual
remains above \(\epsilon\) at every iteration.

\begin{lemma}
\label{lem:rot-trajectory}
Under the conditions of Lemma~\ref{lem:rot-estimates}, the iterates of
\(\mathrm{SPLALM}_{\theta,\alpha,\beta}\) on Example~\ref{ex:rotation}
satisfy, for all \(k\ge0\),
\[
x^k-x_c=\zeta^kX+(\mathbb T^k\br)_x,
\qquad
\lambda^k-\lambda_c=\zeta^k\Lambda+(\mathbb T^k\br)_\lambda,
\]
and \(R_k\ge\frac76\epsilon\). Hence
\(T^{\mathrm{KKT}}_\epsilon(\theta,\alpha,\beta;I)=+\infty\).
\end{lemma}

\begin{proof}
We prove by induction that, for all \(k\ge0\),
\[
\bw^k-\bw_c=\zeta^k\bv+\mathbb T^k\br.
\]
\emph{Base case \(k=0\).} By Lemma~\ref{lem:rot-spectral}~(iii),
\(\bv+\br=-\bw_c\). Since \(\bw^0=\zero\), the identity holds for \(k=0\).

\emph{Inductive step.} Assume that the identity holds for some \(k\ge0\).
Since \(|\zeta|=1\), Lemma~\ref{lem:rot-estimates}~(ii) gives
\[
\frac34X\le\|x^k-x_c\|\le\frac54X=R_*.
\]
Thus, \(x^k\) lies in the quadratic region. Moreover, the induction
hypothesis and Lemma~\ref{lem:rot-spectral}~(iii) imply
\(\|\lambda^k-\lambda_c\|\le\|\Lambda\|+\|(\mathbb T^k\br)_\lambda\|\)
and \(\|\lambda_c\|\le\|\Lambda\|+\|\br_\lambda\|\).
Lemma~\ref{lem:rot-estimates}~(ii) and the parameter bounds in
Lemma~\ref{lem:rot-parameter-bounds} therefore yield
\[
\|\lambda^k+\rho\calA(x^k)\|
\le2\|\Lambda\|+2s\Bigl(\frac{\gamma}{\eta}+\frac54\Bigr)bX+\frac54\rho bX
\le\frac13\Bigl(2+\frac14\Bigr)=\frac34.
\]
Hence \((x^k,\lambda^k)\in\mathcal R_\rho\). By
Proposition~\ref{prop:reduction}, the primal update is given by the first
row of~\eqref{eq:rot-matrix}. Using the induction hypothesis and
\(\mathbb T\bv=\zeta\bv\), we obtain
\[
x^{k+1}-x_c=\zeta^{k+1}X+(\mathbb T^{k+1}\br)_x.
\]
Lemma~\ref{lem:rot-estimates}~(ii) then gives
\(\frac34X\le\|x^{k+1}-x_c\|\le\frac54X\).
The same estimate as above, with \(x^{k+1}\) in place of \(x^k\), yields
\[
\|\lambda^k+\rho\calA(x^{k+1})\|\le\frac34.
\]
Thus, \((x^{k+1},\lambda^k)\in\mathcal R_\rho\) as well, so
Proposition~\ref{prop:reduction} also applies to the dual update.
Together with the center updates, this gives
\[
\bw^{k+1}-\bw_c
=\mathbb T(\bw^k-\bw_c)
=\zeta^{k+1}\bv+\mathbb T^{k+1}\br.
\]
This completes the induction. In particular, both proximal outputs vanish
at every iteration, and all primal iterates lie in the prescribed annulus.

Finally, we bound the residual. By Proposition~\ref{prop:embedding},
\[
R_k=\max\bigl\{\|{-a}(x^{k+1}-x_c)+b(\lambda^k-\lambda_c)\|,\
b\|x^{k+1}-x_c\|\bigr\}.
\]
The second term is at least \(\frac34bX\). For the first term,
Lemma~\ref{lem:rot-estimates}~(ii) and Lemma~\ref{lem:rot-parameter-bounds} give
\[
\|{-a}(x^{k+1}-x_c)+b(\lambda^k-\lambda_c)\|
\ge b\|\Lambda\|-s\Bigl(\frac{\gamma}{\eta}+\frac54\Bigr)b^2X-\frac54aX
\ge\frac{27}{32}b\|\Lambda\|-s\Bigl(\frac{\gamma}{\eta}+\frac52\Bigr)b^2X .
\]
If \(\|\Lambda\|\ge X\), Lemma~\ref{lem:rot-parameter-bounds} bounds the
right-hand side below by \(\frac{23}{32}b\|\Lambda\|\). If
\(X\ge\|\Lambda\|\), we use the second term instead. In both cases,
\(R_k\ge\frac b2\max\{X,\|\Lambda\|\}=\frac b6\ge\frac76\epsilon\).
\end{proof}
 


\subsection{Large primal center stepsize \texorpdfstring{$\alpha$}{alpha}, and \texorpdfstring{$\beta=0$}{beta=0}}
\label{subsec:lb-huber}

In this subsection, we prove the lower bound for \(\alpha>\alpha_\epsilon\)
and \(\beta=0\). The instance below keeps the objective of
Example~\ref{ex:plateau} and changes the inner mapping and the outer
function. On this instance, \(\mathrm{SPLALM}_{\theta,\alpha,0}\), initialized
at zero, never produces an \(\epsilon\)-KKT point, for every
\(\alpha\in[0,1]\).

\begin{example}
\label{ex:huber}
Let \(\epsilon\) satisfy the condition in Example~\ref{ex:plateau}, and let
\(f\) be the function defined there. Set
\[
q_h:=2s+1,
\qquad
a_0:=\frac{4\epsilon(\rho+3s+1)}{q_h}.
\]
Define \(\calA\equiv a_0\) and \(h:\RR\to\RR\) by
\[
h(y)=
\begin{cases}
\dfrac{q_h}{2}y^2, & |y|\le q_h^{-1},\\[4pt]
|y|-\dfrac{1}{2q_h}, & |y|>q_h^{-1}.
\end{cases}
\]
Then \(h\) is convex and \(1\)-Lipschitz, and \(\calA\) satisfies
\(L_{\calA}=G_{\calA}=0\). Since \(F=f+h(a_0)\) differs from \(f\) by a
constant, Remark~\ref{rem:plateau} gives
\(F(0)-\inf F=f(0)-\inf f=\Delta\). Together with the properties of \(f\)
established in Example~\ref{ex:plateau}, this gives
\((f,\calA,h)\in\mathcal I(\mathbf C_0,\Delta)\).
\end{example}

When \(\beta=0\), the dual center stays at \(\mu^k=\mu^0=0\). The dual
perturbation term \(-\frac{1}{2s}\|\lambda\|^2\) in~\eqref{eq:potential}
therefore pulls \(\lambda^k\) toward zero at every iteration. On this
instance, the output \(y^{k+1}\) equals \(a_0\) only when
\(\lambda^k=\nabla h(a_0)=q_ha_0\), whereas the dual iterates converge to a
smaller value. The difference keeps the feasibility residual
\(\|\calA(x^{k+1})-y^{k+1}\|\) above \(\epsilon\) at every iteration. Since
\(J_{\calA}=0\), the multiplier does not enter the primal update, and the
dual iterates have the closed form given in the next lemma.

\begin{lemma}
\label{lem:huber}
Let \(\alpha\in[0,1]\) and \(\beta=0\). Suppose that \(\epsilon\) satisfies
the condition in Example~\ref{ex:plateau} and \(\epsilon(\rho+s)\le c_1\),
with \(c_1\) as specified in Theorem~\ref{thm:lower-bound}. Run
\(\mathrm{SPLALM}_{\theta,\alpha,0}\) on Example~\ref{ex:huber}, and set
\[
\varphi:=1-\frac{\tau}{s}-\frac{\tau}{\rho+q_h},
\qquad
\bar\lambda:=\frac{s\,q_ha_0}{\rho+q_h+s}.
\]
Then the following hold.
\begin{enumerate}[label=(\roman*)]
\item For all \(k\ge0\),
  \[
  \mu^k=0,\qquad
  \lambda^k=\bar\lambda\bigl(1-\varphi^k\bigr),\qquad
  y^{k+1}=\frac{\rho a_0+\lambda^k}{\rho+q_h}.
  \]
\item For all \(k\ge0\),
  \(R_k\ge\|\calA(x^{k+1})-y^{k+1}\|\ge3\epsilon\). Hence, with \(I\) the
  instance of Example~\ref{ex:huber},
  \(T^{\mathrm{KKT}}_\epsilon(\theta,\alpha,0;I)=+\infty\).
\end{enumerate}
\end{lemma}

\begin{proof}
We first collect three bounds on the parameters. Since \(0<\tau\le s\) and
\(\rho+q_h>2s\),
\begin{equation}
\label{eq:huber-ratio}
-\frac12<-\frac{\tau}{\rho+q_h}\le\varphi<1,
\qquad
\frac{s}{\rho+q_h}<\frac12 .
\end{equation}
Hence \(\varphi^k\ge-\frac12\) and \(0\le1-\varphi^k\le2\) for all
\(k\ge0\). Since \(2s\le\rho+q_h+s\), it follows that
\(0\le\bar\lambda(1-\varphi^k)\le2\bar\lambda\le q_ha_0\). Finally,
\(s\ge\tau\) gives \(1\le(\rho+s)/\tau\), so the third term in the
definition of \(c_1\) yields
\[
q_ha_0=4\epsilon(\rho+3s+1)
\le4\epsilon\bigl(3+\tau^{-1}\bigr)(\rho+s)
\le4c_1\bigl(3+\tau^{-1}\bigr)
\le\frac14 .
\]

(i) Since \(\beta=0\) and \(\mu^0=0\), the dual center update gives
\(\mu^k=0\) for all \(k\). Next, fix \(x\in\RR\) and
\(\lambda\in[0,q_ha_0]\), and set
\(y_\lambda:=(\rho a_0+\lambda)/(\rho+q_h)\). Then
\(0\le y_\lambda\le a_0\le q_h^{-1}\), where the last inequality uses
\(q_ha_0\le\frac14\). On \([-q_h^{-1},q_h^{-1}]\), the function \(h\) is
differentiable with \(\nabla h(y)=q_hy\), and \(y_\lambda\) satisfies
\[
q_hy_\lambda-\lambda-\rho\bigl(a_0-y_\lambda\bigr)=0,
\]
which is the optimality condition of the strongly convex problem
in~\eqref{eq:reduced-al}. Hence \(y(x,\lambda)=y_\lambda\), and
\[
\calA(x)-y(x,\lambda)=\frac{q_ha_0-\lambda}{\rho+q_h}.
\]
By Danskin's theorem,
\(\nabla_\lambda\calD_\rho(x,\lambda)=\calA(x)-y(x,\lambda)\).

We now prove the formula for \(\lambda^k\) by induction. It holds for
\(k=0\) because \(\lambda^0=0\). Assume that it holds for some \(k\ge0\).
Then \(\lambda^k\in[0,q_ha_0]\), so
\(y^{k+1}=y(x^{k+1},\lambda^k)=y_{\lambda^k}\). Together with
\(\mu^k=0\), the dual update in~\eqref{eq:splalm-updates} reads
\[
\lambda^{k+1}
=\lambda^k+\tau\Bigl[\frac{q_ha_0-\lambda^k}{\rho+q_h}
-\frac{\lambda^k}{s}\Bigr]
=\varphi\lambda^k+\frac{\tau q_ha_0}{\rho+q_h}.
\]
A direct computation shows that
\(\bar\lambda=\varphi\bar\lambda+\tau q_ha_0/(\rho+q_h)\). Subtracting the
two identities gives
\(\lambda^{k+1}-\bar\lambda=\varphi(\lambda^k-\bar\lambda)
=-\bar\lambda\varphi^{k+1}\), which completes the induction.

(ii) Since \(J_{\calA}=0\), \eqref{eq:kkt-residual} gives
\(R_k=\max\{\|\nabla f(x^{k+1})\|,\|\calA(x^{k+1})-y^{k+1}\|\}\). By (i)
and \(q_ha_0-\bar\lambda=q_ha_0(\rho+q_h)/(\rho+q_h+s)\),
\[
\calA(x^{k+1})-y^{k+1}
=\frac{q_ha_0-\lambda^k}{\rho+q_h}
=\frac{q_ha_0}{\rho+q_h+s}\Bigl(1+\frac{s}{\rho+q_h}\varphi^k\Bigr)
=4\epsilon\Bigl(1+\frac{s}{\rho+q_h}\varphi^k\Bigr),
\]
where the last equality uses \(\rho+q_h+s=\rho+3s+1\).
By~\eqref{eq:huber-ratio}, \(\frac{s}{\rho+q_h}\varphi^k>-\frac14\), so
\(\|\calA(x^{k+1})-y^{k+1}\|>3\epsilon\). Thus no output is an
\(\epsilon\)-KKT point, and
\(T^{\mathrm{KKT}}_\epsilon(\theta,\alpha,0;I)=+\infty\).
\end{proof}

\subsection{Proof of Theorem~\ref{thm:lower-bound}}
\label{subsec:lb-combine}

We now combine the three instances. Let \(c_1\) and \(c_2\) be given
by~\eqref{eq:rot-thresholds}. The constant \(c_1\) depends only on
\((\gamma,\eta,\tau)\), and \(c_2\) depends only on \((\gamma,\tau,\Delta)\),
as required in Theorem~\ref{thm:lower-bound}.

\begin{proof}[Proof of Theorem~\ref{thm:lower-bound}]
Fix \(\epsilon\in(0,c_2)\), \(\rho>0\), and \(s\ge\tau\) with
\(\epsilon(\rho+s)\le c_1\), and fix \((\alpha,\beta)\in[0,1]^2\). The
definition of \(c_2\) gives \(\epsilon\le\frac1{64}\le\frac12\),
\(\epsilon\le\sqrt\Delta/2\), and
\(\epsilon^2\le\Delta/(64\gamma)\le\Delta/(16\gamma)\). Hence the conditions
of Example~\ref{ex:plateau}, Lemma~\ref{lem:plateau}, and
Example~\ref{ex:huber} on \(\epsilon\) hold. We treat the three parts of the
partition of \([0,1]^2\) separately.

\emph{Case 1: \(\alpha\le\alpha_\epsilon\).} Let \(I\) be the instance of
Example~\ref{ex:plateau}, which belongs to \(\mathcal I(\mathbf C_0,\Delta)\).
If \(\alpha=0\), Lemma~\ref{lem:plateau}~(iii) gives
\(T^{\mathrm{KKT}}_\epsilon(\theta,\alpha,\beta;I)=+\infty\). If
\(0<\alpha\le\alpha_\epsilon\), Lemma~\ref{lem:plateau}~(iii) and the
definition of \(\alpha_\epsilon\) give
\[
T^{\mathrm{KKT}}_\epsilon(\theta,\alpha,\beta;I)
>\frac{\Delta}{16\gamma\alpha\epsilon^2}
\ge\frac{\Delta}{16\gamma\alpha_\epsilon\epsilon^2}
=\frac{\Delta\,\epsilon^{-4}}{16384\,\gamma\max\{1,4\tau,4\gamma\tau\}}.
\]
Both bounds hold for every \(\beta\in[0,1]\).

\emph{Case 2: \(\alpha>\alpha_\epsilon\) and \(\beta>0\).} The conditions of
Lemma~\ref{lem:rot-estimates} hold. Let \(I\) be the instance of
Example~\ref{ex:rotation}. By Lemma~\ref{lem:rot-estimates}~(i),
\(I\in\mathcal I(\mathbf C_0,\Delta)\), and Lemma~\ref{lem:rot-trajectory}
gives \(T^{\mathrm{KKT}}_\epsilon(\theta,\alpha,\beta;I)=+\infty\).

\emph{Case 3: \(\alpha>\alpha_\epsilon\) and \(\beta=0\).} The conditions of
Lemma~\ref{lem:huber} hold. Let \(I\) be the instance of
Example~\ref{ex:huber}, which belongs to \(\mathcal I(\mathbf C_0,\Delta)\).
Lemma~\ref{lem:huber}~(ii) gives
\(T^{\mathrm{KKT}}_\epsilon(\theta,\alpha,0;I)=+\infty\).

The three cases cover \([0,1]^2\). Hence, for every
\((\alpha,\beta)\in[0,1]^2\), some \(I\in\mathcal I(\mathbf C_0,\Delta)\)
satisfies
\[
T^{\mathrm{KKT}}_\epsilon(\theta,\alpha,\beta;I)
\ge\frac{\Delta\,\epsilon^{-4}}{16384\,\gamma\max\{1,4\tau,4\gamma\tau\}}.
\]
Taking the supremum over \(I\) and then the infimum over \((\alpha,\beta)\)
in Definition~\ref{def:complexity-measure} gives the bound on
\(T^{\mathrm{KKT}}_\epsilon(\calC_\theta,\mathcal I(\mathbf C_0,\Delta))\).
\end{proof}

In Cases~2 and~3, the method never produces an \(\epsilon\)-KKT output.
The finite bound comes only from Case~1, and it determines the order
\(\epsilon^{-4}\) in Theorem~\ref{thm:lower-bound}.
